\subsection{But et objectifs}

Le laboratoire réalisé est constitué de deux manipulations concernant les ondes. Ci-dessous
figure la liste des objectifs à réaliser, provenant des descriptifs fournis en classe \citeit{descriptif_tpe11} \citeit{descriptif_tpe12}.

\textbf{Manipulation 1 : Ondes centimétriques -- Micro-ondes}

\begin{itemize}
  \item Objectif 1 : Propagation rectiligne -- montrer que le signal est maximal lorsque les cornets sont alignés.
  \item Objectif 2 : Capacité de pénétration -- vérifier que les micro-ondes traversent une plaque de bois.
  \item Objectif 3 : Blindage et absorption -- comparer la transmission sans plaque, avec le bois et avec l'aluminium.
  \item Objectif 4 : Réflexion -- vérifier que l'angle de réflexion est égal à l'angle d'incidence.
  \item Objectif 5 : Ondes stationnaires -- déterminer la longueur d'onde et la fréquence de l'émetteur.
  \newpage
  \item Objectif 6 : Réfraction -- mesurer l'indice de réfraction de la paraffine à l'aide d'un prisme.
  \item Objectif 7 : Principe de Huygens -- montrer qu'une fente renvoie du signal dans la zone d'ombre.
  \item Objectif 8 : Diffraction -- observer le signal derrière une plaque couvrante, dans la zone d'ombre.
  \item Objectif 9 : Interférence -- repérer les maxima et minima produits par une double fente.
  \item Objectif 10 : Polarisation -- comparer le signal selon l'orientation de la grille de polarisation.
\end{itemize}

\textbf{Manipulation 2 : Trombone de Koenig -- Vitesse du son dans les gaz}

\begin{itemize}
  \item Objectif 1 : Vitesse théorique -- calculer la vitesse attendue dans l'air et le CO$_2$.
  \item Objectif 2 : Mesure dans l'air -- déterminer la vitesse du son par régression linéaire sur plusieurs fréquences.
  \item Objectif 3 : Mesure dans CO$_2$ -- déterminer la vitesse du son par régression linéaire sur plusieurs fréquences.
  \item Objectif 4 : Comparaison -- vérifier que la vitesse ne dépend pas de la fréquence et suit la loi théorique.
\end{itemize}
